Paired-sample $t$-tests between synthetic and real images revealed no statistically significant differences in normalized regional uptake across the principal brain regions shown in Fig.~\ref{fig:4}A, B for the overall cohorts of both tracers ($p>0.05$). Across all 98 ROIs, no overall comparisons remained significant after Benjamini--Hochberg false discovery rate (FDR) correction (all $q>0.05$). Among the diagnosis-stratified paired comparisons, only the right pericalcarine cortex (ROI 2021) in the $^{18}$F-FDG MCI subgroup showed a significant mean paired difference after FDR correction ($q=0.0481$). The median concordance correlation coefficients (CCC) across the 98 ROIs were 0.9745 for $^{18}$F-AV-45 and 0.9450 for $^{18}$F-FDG, with corresponding median mean absolute errors (MAE) of 0.0160 and 0.0159 (normalized intensity units). For the 16 representative tracer--ROI combinations in Fig.~\ref{fig:4}A, B, descriptive regression slopes ranged from 0.9794 to 1.0127 (median, 1.0007) for $^{18}$F-AV-45 and from 0.9844 to 1.0101 (median, 1.0001) for $^{18}$F-FDG, without a consistent direction of deviation from unity. These results indicate close agreement in normalized regional uptake between synthetic and real PET, with detailed statistics for the representative ROIs provided in Supplementary Table~\ref{tab:regional_agreement}. Non-significant mean differences do not establish formal equivalence or clinical interchangeability. These historical regional-CSV results are separate from the reconstructed Full generation comparisons.
